\documentclass[../../../include/open-logic-section]{subfiles}

\begin{document}

\olfileid{sth}{cardinals}{cp}
\olsection{કૅન્ટરનો સિદ્ધાંત}

\olref[ordinals][vn]{sec} યાદ કરો. ત્યાં આપણે સુક્રમિત
ગણોની ચર્ચા કરી અને સુક્રમોના અવેજ તરીકે કામ કરે એવી
વસ્તુઓ હોય તો સારું એમ સૂચવ્યું હતું. આ વિચારથી
આપણે ક્રમવાચક સંખ્યાઓ રજૂ કરી; પછી
\olref[ordinals][ordtype]{ordtypesworklikeyouwant}માં
બતાવ્યું કે તેઓ અપેક્ષા પ્રમાણે વર્તે છે, એટલે કે:
\[
	\ordtype{A, <} = \ordtype{B, \lessdot} 
	\text{ iff } \tuple{A, <} \isomorphic \tuple{B, \lessdot}.
\]
હવે તેનાથી પણ પહેલાંના \olref[sfr][siz][equ]{sec}ને યાદ કરો.
ત્યાં આપણે અનૌપચારિક રીતે કામ કરતાં ગણના ``કદ''નો
ખ્યાલ રજૂ કર્યો હતો. ખાસ કરીને, બે ગણો સામ્ય છે,
એટલે $\cardeq{A}{B}$, જો અને માત્ર જો કોઈ
!!a{bijection} $f \colon A \to B$ હોય, એમ કહ્યું હતું.
આ ખ્યાલ સુક્રમ કરતાં મૂળભૂત રીતે {સરળ} છે:
આપણને ફક્ત !!{bijection}sમાં રસ છે; ક્રમ-એકરૂપતાઓની
જેમ એ !!{bijection}s કોઈ રચના ``જાળવે'' છે કે નહીં તેમાં નહીં.

આ બધી વાતમાંથી એક સ્વાભાવિક વિચાર મળે છે.
જેમ સુક્રમોનું માપ દર્શાવવા માટે આપણે કેટલીક વસ્તુઓ,
એટલે \emph{ક્રમવાચક સંખ્યાઓ}, રજૂ કરી, તેમ કદનું માપ
દર્શાવવા માટે કેટલીક વસ્તુઓ, એટલે \emph{ગણાંકો},
રજૂ કરી શકીએ. આ પ્રકરણનો હેતુ એ જ છે.

આ ગણાંકો શું હશે તે કહીએ તે પહેલાં તેમણે સંતોષવો
જોઈએ એવો સિદ્ધાંત નક્કી કરીએ. ગણ $X$નો ગણાંક
દર્શાવવા $\card{X}$ લખીએ, તો આપણે ઇચ્છીએ કે:
\[
	\card{A} = \card{B} \text{ iff } \cardeq{A}{B}.
\]
આને \emph{કૅન્ટરનો} સિદ્ધાંત કહીશું, કારણ કે કદાચ
કૅન્ટરે જ સૌપ્રથમ આ વિચાર સ્પષ્ટ રીતે મનમાં રાખ્યો હતો.
(\olref[hp]{sec}માં \emph{હ્યુમના} સિદ્ધાંત સાથેનો તેનો
સંબંધ વધુ જોઈશું.) તેથી અમારો હેતુ દરેક $X$ માટે
$\card{X}$ એવી રીતે વ્યાખ્યાયિત કરવાનો છે કે કૅન્ટરનો
સિદ્ધાંત મળે.

\end{document}
